%auto-ignore
\section{서론}


언어모델 사전학습(language model pre-training)은 많은 자연어 처리(natural language processing) 과제의 성능을 끌어올리는 데 효과적임이 보여져 왔다~\cite{dai-le:2015:_semi, peters-etal:2018:_deep, radford-etal:2018, howard-ruder:2018}. 여기에는 문장 사이의 관계를 통째로 분석해 예측하는 자연어 추론(natural language inference)~\cite{bowman-etal:2015, williams-nangia-bowman:2018}과 패러프레이즈(paraphrasing)~\cite{dolan-brockett:2005:_autom} 같은 문장 수준 과제(sentence-level task), 그리고 모델이 토큰 수준의 미세한 출력을 만들어야 하는 개체명 인식(named entity recognition)과 질의응답 같은 토큰 수준 과제(token-level task)~\cite{tjong-de:2003, rajpurkar-etal:2016:_squad}가 포함된다.


사전학습된 언어 표현을 다운스트림 과제(downstream task)에 적용하는 데에는 두 가지 기존 전략이 있다 — \emph{피처 기반(feature-based)} 과 \emph{파인튜닝(fine-tuning)}. ELMo~\cite{peters-etal:2018:_deep} 같은 피처 기반 접근은, 사전학습 표현을 추가 피처(feature)로 포함하는 과제별 아키텍처(task-specific architecture)를 사용한다. Generative Pre-trained Transformer(OpenAI GPT)~\cite{radford-etal:2018} 같은 파인튜닝 접근은 과제별 파라미터(task-specific parameter)를 최소한만 추가하고, 사전학습된 \emph{모든} 파라미터를 단순히 파인튜닝하는 방식으로 다운스트림 과제에서 학습한다. 두 접근은 사전학습 단계에서 동일한 목적함수를 공유하며, 일반적인 언어 표현을 학습하기 위해 단방향(unidirectional) 언어모델을 사용한다.

저자들의 주장은, 현재 기법들이 사전학습 표현의 잠재력을 제약하고 있다는 것이다 — 특히 파인튜닝 기반 접근에서. 가장 큰 한계는 표준 언어모델이 단방향(unidirectional)이라는 점, 그리고 이로 인해 사전학습에 사용할 수 있는 아키텍처 선택이 제한된다는 점이다. 예를 들어 OpenAI GPT에서는 좌→우(left-to-right) 아키텍처를 사용하므로, 트랜스포머(Transformer)~\cite{vaswani-etal:2017:_atten}의 셀프 어텐션(self-attention) 층에서 모든 토큰(token)이 자기 이전 토큰에만 어텐션(attention)할 수 있다. 이런 제약은 문장 수준 과제에서 차선이며, 질의응답처럼 양쪽 맥락을 모두 반영해야 하는 토큰 수준 과제에 파인튜닝 기반 접근을 적용할 때는 매우 해로울 수 있다.

이 논문에서 저자들은 \bert: \textbf{B}idirectional \textbf{E}ncoder \textbf{R}epresentations from \textbf{T}ransformers를 제안하여 파인튜닝 기반 접근을 개선한다. \bert는 ``마스킹 언어모델(masked language model, MLM)''이라는 사전학습 목적함수(pre-training objective)를 사용해 앞서 언급한 단방향성 제약을 완화한다. 이 목적함수는 Cloze 과제~\cite{taylor:1953:_cloze}에서 영감을 얻었다. 마스킹 언어모델은 입력 토큰의 일부를 무작위로 마스킹(masking)하고, 마스킹된 단어의 원본 어휘 id를 그 단어의 맥락만으로 예측하는 것을 목적으로 한다. 좌→우 언어모델 사전학습과 달리, MLM 목적함수는 표현이 좌·우 맥락을 융합할 수 있게 해 주어, 깊은 양방향 트랜스포머(deep bidirectional Transformer)를 사전학습할 수 있게 한다. 마스킹 언어모델 외에, 저자들은 텍스트 쌍 표현을 함께 사전학습하는 ``다음 문장 예측(next sentence prediction, NSP)'' 과제도 사용한다. 본 논문의 기여는 다음과 같다.
\begin{itemize}[leftmargin=1em]
  \item 언어 표현 학습에서 양방향 사전학습의 중요성을 입증한다. 사전학습에 단방향 언어모델을 사용한 \citet{radford-etal:2018}과 달리, \bert는 마스킹 언어모델을 사용해 사전학습된 깊은 양방향 표현을 가능하게 한다. 또한 독립적으로 학습된 좌→우와 우→좌 LM의 얕은 결합(concatenation)을 사용한 \citet{peters-etal:2018:_deep}과도 대비된다.
  \item 사전학습된 표현이 무겁게 엔지니어링된 과제별 아키텍처의 필요성을 줄여 준다는 것을 보인다. \bert는 다수의 문장 수준 \emph{및} 토큰 수준 과제에서 최첨단 성능을 달성한 최초의 파인튜닝 기반 표현 모델이며, 많은 과제별 아키텍처를 능가한다.
  \item \bert는 11개 NLP 과제에서 최첨단을 갱신한다. 코드와 사전학습된 모델은 \url{https://github.com/google-research/bert}에서 받을 수 있다.
\end{itemize}
